\section{Conclusion}
We asked whether a language model's internal state shows that newly retrieved evidence has become sufficient, and whether the model will use it. Both are readable from the residual stream at the answer position before generation, in three model families and under splits that keep documents, answers and relations disjoint. For sufficiency the internal state is a far better judge than the model's own output, which declares complete evidence insufficient a third of the time, and somewhat better than a text classifier given the whole context. For uptake, text alone says almost nothing, the model's decision to answer says much, and the internal state adds to it. The directions that carry these readouts are not the ones the model acts on; the only causal direction we found is a gate on abstention. The readouts transfer across datasets, fully to HotpotQA and partially to other domains and to retrieved evidence. Beyond what they say about the model, these readouts have practical uses: read before generation, the sufficiency score can decide when an adaptive retrieval system should stop fetching evidence, and the uptake score whether to trust the answer it then gives (\cref{app:adaptive}). We release the trajectory construction, the probing and intervention code, and the benchmark.
